%%%%%%%%%%  scan idx 28  |  folio 15  %%%%%%%%%%
la
premi\`ere fl\^eche verticale est la m\^eme que dans (5.3), et la deuxi\`eme est l'homomorphisme de Gysin en cohomologie rationelle.

\subsection*{6.}
Dans tout ceci, il n'a \'et\'e gu\`ere question que de l'invariant contravariant $K^\bullet(X)$, \`a l'exclusion de l'invariant covariant $K_\bullet(X)$ dont l'existence a \'et\'e signal\'e dans 3.1. Si on consid\`ere $K^\bullet(X)$ comme donnant un analogue de l'anneau de \underline{cohomologie} enti\`ere d'un espace topologique, on doit consid\'erer $K_\bullet(X)$ comme donnant un analogue de l'\underline{homologie} enti\`ere, la structure de Module de $K_\bullet(X)$ sur $K^\bullet(X)$ correspondant alors \`a l'op\'eration de cap-produit\footnote{Cette analogie n'est d'ailleurs pas purement verbale, et peut se pr\'eciser, apr\`es passage aux gradu\'es $\Gr^\bullet(X)$, $\Gr_\bullet(X)$ associ\'es, par une relation de compatibilit\'e explicite entre la structure de module sur $\Gr_\bullet(X)$ et l'op\'eration «~cap~».}. Le fait que lorsque $X$ est r\'egulier, l'homomorphisme naturel $K^\bullet(X)\to K_\bullet(X)$ soit un isomorphisme, doit \^etre regard\'e comme l'analogue du th\'eor\`eme de dualit\'e de Poincar\'e sur les vari\'et\'es topologiques orient\'ees, affirmant que le cap-produit avec la classe d'homologie fondamentale d\'efinit un isomorphisme de la cohomologie sur l'homologie. Ces analogies sugg\`erent d'ailleurs l'opportunit\'e de d\'efinir \'egalement, pour un poly\`edre fini $X$ mettons, un invariant $K_\bullet(X)$ de nature covariante, dont les relations \`a l'homologie enti\`ere (suite spectrale, homomorphisme de Chern\dots) seraient analogues \`a celles existant entre $K^\bullet(X)$ et la cohomologie enti\`ere~; et il ne semble pas non plus exclu qu'il puisse exister, en termes de ces groupes $K_\bullet(X)$, des variantes du th\'eor\`eme de Riemann--Roch diff\'erentiable, valable pour des espaces plus g\'en\'eraux que des vari\'et\'es diff\'erentiables compactes.

Revenant au cas o\`u $X$ est un sch\'ema noeth\'erien quelconque, il convient de munir $K_\bullet(X)$ d'une filtration croissante, $\operatorname{Filt}_iK_\bullet(X)$ \'etant engendr\'e par des classes $\cl_\bullet(F)$, avec $F$ faisceau coh\'erent sur $X$ tel que $\dim \operatorname{supp}.F\le i$. On prouvera (Exp. X) la compatibilit\'e de cette filtration avec la structure de module et la filtration de $K^\bullet(X)$, i.e. la formule
\begin{equation}\tag{6.1}
\operatorname{Filt}^iK^\bullet(X)\cdot \operatorname{Filt}_j(K_\bullet(X)\subset \operatorname{Filt}_{j-i}K_\bullet(X),
\end{equation}
ce qui permet de munir le gradu\'e associ\'e au module filtr\'e $K_\bullet(X)$, not\'e
\[
\Gr_\bullet(X),
\]
d'une structure de module sur $\Gr^\bullet(X)$. Ses \'el\'ements peuvent d'ailleurs s'interpr\'eter comme des \uline{classes de cycles alg\'ebriques} sur $X$, (pour une relation d'\'equivalence moins fine que l'\'equivalence rationelle, lorsque $X$ est de type fini sur un corps de base $K$). Si
